\chapter{Análisis de resultados}
\label{chap:resultados}

En este capítulo se presentan y analizan los resultados obtenidos de las pruebas realizadas con los compiladores \emph{GCC}, \emph{Clang} y \emph{Rustc}, evaluando el impacto de diferentes niveles de optimización y medidas de seguridad en cuatro métricas: tiempo de ejecución, uso de memoria, tamaño del archivo generado y porcentaje de funciones fortificadas. Los experimentos se han llevado a cabo en un entorno controlado utilizando programas específicamente diseñados para aislar los efectos de las combinaciones evaluadas. El análisis comparativo revela interesantes compensaciones entre rendimiento y seguridad, destacando las fortalezas particulares de cada compilador en diferentes escenarios.

\section{Entorno de experimentación}

Las pruebas se han llevado a cabo en un entorno controlado, utilizando
un único ordenador para asegurar la coherencia de los resultados. A
continuación se detallan las características principales del ordenador
donde se han realizado las pruebas:

\begin{itemize}
\item
  \textbf{Procesador}: Intel Core i7-10750H (6 núcleos físicos, 12
  hilos) con frecuencia base de 2.60\,GHz y turbo de hasta 5.00\,GHz
\item
  \textbf{Memoria RAM}: 16\,GB
\item
  \textbf{Tarjeta gráfica}: NVIDIA GeForce RTX 3060 Mobile / Max-Q
\item
  \textbf{Sistema operativo}: Kali GNU/Linux Rolling (x86\_64)
\item
  \textbf{Kernel}: Linux 6.12.25
\end{itemize}

Para realizar las pruebas se emplearon las siguientes versiones de
compiladores:

\begin{itemize}
\item
  \emph{GCC}: versión 14.2.0 (Debian 14.2.0-19)
\item
  \emph{Clang}: versión 19.1.7 (Debian)
\item
  \emph{Rustc}: versión 1.89.0-nightly
\end{itemize}

La elección de la versión nightly de \emph{Rustc} es debido a que las
opciones avanzadas de seguridad, como los sanitizadores
(\texttt{-Z\ sanitizer=...}) y los protectores de pila
(\texttt{-Z\ stack-protector=...}) no están disponibles en la rama
estable.

\section{Descripción de conjuntos de datos}

Para realizar las pruebas se han utilizado dos programas escritos en C++
y sus equivalentes en Rust. Estos programas han sido diseñados para
activar de forma controlada las protecciones de seguridad evaluadas, y
para generar un nivel de trabajo suficiente como para observar el efecto
de las optimizaciones.


Programa \underline{\textbf{suma.cpp}} / \underline{\textbf{suma.rs}}.
Este programa realiza una operación intensiva sobre un vector de tamaño
considerable (1 millón de posiciones), inicializado con un valor
constante.

En ambos casos, el programa permite observar el impacto de la optimización
sobre bucles y uso de estructuras (\texttt{std::vector} en C++, \texttt{Vec} en Rust), 
los efectos sobre el uso de memoria y el tamaño del binario, y la activación de protecciones
básicas sin provocar errores.\\

Programa \underline{\textbf{final.cpp}} / \underline{\textbf{final.rs}}.
Este programa está específicamente diseñado para mostrar
vulnerabilidades que pueden ser detectadas por mecanismos como
\emph{AddressSanitizer}, \emph{MemorySanitizer} o
\emph{stack\ protector}. El programa realiza una copia insegura de una cadena de 
entrada fija con tamaño de 47 bytes a un búfer pequeño de 32 bytes, 
lo que produce un desbordamiento de memoria en la pila si 
no se aplican medidas de protección.

Las características más releventase son las siguientes: en C++, el uso de 
\texttt{strcpy} (función insegura) junto con un búfer de 32 bytes 
activa múltiples protecciones. La versión de Rust simula un comportamiento 
similar mediante copia manual sin validación de límites. En ambos casos, se 
incluye uso de memoria dinámica para ampliar el análisis.

\bigskip

Ambos programas son breves y están pensados para aislar el efecto de las
medidas de seguridad y las optimizaciones del compilador, evitando
interferencias provocadas por lógica compleja o dependencias externas.
Esto facilita la interpretación de los resultados, especialmente en
gráficos donde se comparan directamente métricas como el tiempo de
ejecución, el uso de memoria o el porcentaje de funciones fortificadas.


Las versiones fueron compiladas sin modificar para cada compilador
evaluado: \emph{GCC}, \emph{Clang} y \emph{Rustc}. No dependen de
bibliotecas externas más allá de la biblioteca estándar, lo que garantiza la
reproducibilidad del experimento en distintos entornos.

\section{Descripción de experimentos y resultados}

\subsection{Diseño del experimento}

Para cada programa —uno en C++ y su versión equivalente en Rust— se
generan múltiples binarios combinando una única flag de seguridad,
seleccionada de entre un conjunto representativo, y un nivel de optimización,
específico de cada compilador.

Esto da lugar a un conjunto amplio de combinaciones para cada lenguaje y
compilador, que posteriormente se comparan en términos de tiempo de ejecución, 
uso de memoria, tamaño del binario resultante y porcentaje de funciones fortificadas.
Cada experimento se ha ejecutado 10 veces para calcular valores medios.

\subsection{Combinaciones}

A continuación se resumen las \emph{flags} con las que se realizan las
combinaciones evaluadas para cada compilador. Las tablas \ref{tab:gcc_optimizacion},
\ref{tab:clang_optimizacion} y \ref{tab:rust_optimizacion} muestran las opciones de
optimización para cada compilador, mientras que las opciones de seguridad se encuentran en las Tablas
\ref{tab:gcc_seguridad}, \ref{tab:clang_seguridad} y \ref{tab:rust_seguridad}.


\begin{table}[H]
\centering
\small % Tamaño de fuente más pequeño
\begin{tabularx}{\linewidth}{|l|Y|}
\hline
\textbf{\emph{Flag}} & \textbf{Descripción} \\
\hline
\texttt{-O0} & Sin optimización. Compilación rápida, ejecución lenta. \\
\texttt{-O1} & Optimización básica, reduce algo el tamaño y mejora rendimiento. \\
\texttt{-O2} & Optimización estándar, buena relación entre velocidad y tamaño. \\
\texttt{-O3} & Optimización agresiva, busca máximo rendimiento (puede aumentar el tamaño). \\
\texttt{-Os} & Optimiza para reducir el tamaño del binario. \\
\hline
\end{tabularx}
\caption{Niveles de optimización en \emph{GCC} [15]}
\label{tab:gcc_optimizacion}
\end{table}

\begin{table}[H]
\centering
\footnotesize % Tamaño de fuente aún más pequeño para tabla ancha
\begin{tabularx}{\linewidth}{|l|Y|}
\hline
\textbf{\emph{Flag}} & \textbf{Descripción} \\
\hline
\texttt{(vacío)} & Sin protecciones adicionales. \\
\texttt{-no-pie} & Ejecutable con dirección de carga fija (menos seguro). \\
\texttt{-pie} & Ejecutable con dirección de carga aleatoria (\emph{PIE}). \\
\texttt{-D\_FORTIFY\_SOURCE=2} & Fortificación avanzada de funciones inseguras. \\
\texttt{-D\_FORTIFY\_SOURCE=1} & Fortificación básica. \\
\texttt{-U\_FORTIFY\_SOURCE} & Desactiva la fortificación. \\
\texttt{-z noexecstack} & Evita que la memoria de ejecución tenga permiso de ejecución. \\
\texttt{-z execstack} & Permite ejecutar código en memoria (menos seguro). \\
\texttt{-fno-plt} & Desactiva tabla de enlaces indirectos (ligero refuerzo de seguridad). \\
\texttt{-fsanitize=address} & Detecta accesos inválidos a memoria (\emph{AddressSanitizer}). \\
\texttt{-fsanitize=undefined} & Detecta comportamientos indefinidos del lenguaje. \\
\texttt{-Wl,-z,relro,-z,now} & Protege la tabla de enlaces contra sobrescrituras (\emph{Full RELRO}). \\
\texttt{-fno-stack-protector} & Desactiva protección contra desbordamiento de variables locales. \\
\texttt{-fstack-protector} & Protección básica contra desbordamiento de pila. \\
\texttt{-fstack-protector-strong} & Protección más amplia contra desbordamientos. \\
\texttt{-fstack-protector-all} & Protección total: cualquier función con variables en pila será protegida. \\
\hline
\end{tabularx}
\caption{\emph{Flags} de seguridad en \emph{GCC} [16]}
\label{tab:gcc_seguridad}
\end{table}


\begin{table}[H]
\centering
\small
\begin{tabularx}{\linewidth}{|l|Y|}
\hline
\textbf{\emph{Flag}} & \textbf{Descripción} \\
\hline
\texttt{-O0} & Sin optimización. Compilación rápida, ejecución lenta. \\
\texttt{-O1} & Optimización básica, reduce algo el tamaño y mejora rendimiento. \\
\texttt{-O2} & Optimización estándar, buena relación entre velocidad y tamaño. \\
\texttt{-O3} & Optimización agresiva, busca máximo rendimiento (puede aumentar el tamaño). \\
\texttt{-Os} & Optimiza para reducir el tamaño del binario. \\
\texttt{-Oz} & Optimiza para reducir el tamaño del binario al máximo. \\
\texttt{-Ofast} & Optimiza con máxima agresividad, desactiva comprobaciones de precisión numérica. \\
\hline
\end{tabularx}
\caption{Niveles de optimización en \emph{Clang} [17]}
\label{tab:clang_optimizacion}
\end{table}

\begin{table}[H]
\centering
\footnotesize
\begin{tabularx}{\linewidth}{|l|Y|}
\hline
\textbf{\emph{Flag}} & \textbf{Descripción} \\
\hline
\texttt{(vacío)} & Sin protecciones adicionales. \\
\texttt{-fstack-protector} & Igual que en \emph{GCC}. \\
\texttt{-fstack-protector-strong} & Igual que en \emph{GCC}. \\
\texttt{-fstack-protector-all} & Igual que en \emph{GCC}. \\
\texttt{-fno-stack-protector} & Igual que en \emph{GCC}. \\
\texttt{-pie -f\emph{PIE}} & Activa ejecutable independiente de la dirección de carga (\emph{PIE}). \\
\texttt{-no-pie} & Desactiva \emph{PIE}. \\
\texttt{-D\_FORTIFY\_SOURCE=2} & Nivel de fortificación. \\
\texttt{-D\_FORTIFY\_SOURCE=1} & Nivel de fortificación. \\
\texttt{-D\_FORTIFY\_SOURCE=0} & Desactivación. \\
\texttt{-z noexecstack} & Evita que la memoria de ejecución tenga permiso de ejecución. \\
\texttt{-Wl,-z,execstack} & Permite ejecutar código en memoria (menos seguro). \\
\texttt{-fno-plt} & Desactiva la tabla de enlaces indirectos. \\
\texttt{-fsanitize=address} & Sanitizadores para errores de memoria. \\
\texttt{-fsanitize=undefined} & Sanitizadores para comportamiento indefinido. \\
\texttt{-Wl,-z,relro,-z,now} & Protege estructuras internas al enlazar ficheros (\emph{Full RELRO}). \\
\texttt{-flto} & Activación de optimization al enlazar ficheros (global entre archivos). \\
\texttt{-fsanitize=safe-stack} & Protección específica: separa datos vulnerables de otras variables en pila. \\
\texttt{-fsanitize=memory} & Detecta uso de memoria sin inicializar (\emph{MemorySanitizer}). \\
\hline
\end{tabularx}
\caption{\emph{Flags} de seguridad en \emph{Clang} [17]}
\label{tab:clang_seguridad}
\end{table}


\begin{table}[H]
\centering
\small
\begin{tabularx}{\linewidth}{|l|Y|}
\hline
\textbf{\emph{Flag}} & \textbf{Descripción} \\
\hline
\texttt{-C opt-level=0} & Sin optimización. \\
\texttt{-C opt-level=1} & Optimización básica. \\
\texttt{-C opt-level=2} & Optimización estándar. \\
\texttt{-C opt-level=3} & Optimización máxima. \\
\texttt{-C opt-level=s} & Optimiza para tamaño reducido. \\
\texttt{-C opt-level=z} & Optimiza para tamaño mínimo absoluto. \\
\hline
\end{tabularx}
\caption{Niveles de optimización en \emph{Rustc} [18]}
\label{tab:rust_optimizacion}
\end{table}

\begin{table}[H]
\centering
\footnotesize
\begin{tabularx}{\linewidth}{|l|Y|}
\hline
\textbf{\emph{Flag}} & \textbf{Descripción} \\
\hline
\texttt{(vacío)} & Sin medidas adicionales. \\
\texttt{-C overflow-checks=on} & Comprobación en tiempo de ejecución de desbordamientos aritméticos. \\
\texttt{-C debug-assertions=on} & Comprobaciones de depuración activadas. \\
\texttt{-C panic=abort} & El programa se aborta directamente ante un error grave (más rápido, menos seguro). \\
\texttt{-C lto} & Optimization al enlazar ficheros (convencional). \\
\texttt{-C lto=thin} & Optimization al enlazar ficheros (más ligera). \\
\texttt{-Z sanitizer=address} & Detecta accesos inválidos a memoria (\emph{AddressSanitizer}). \\
\texttt{-Z sanitizer=leak} & Detecta fugas de memoria (\emph{LeakSanitizer}). \\
\texttt{-Z sanitizer=memory} & Detecta uso de memoria sin inicializar (\emph{MemorySanitizer}). \\
\texttt{-C link-arg=-Wl,-z,relro,-z,now} & Protege estructuras internas al enlazar (\emph{Full RELRO}). \\
\texttt{-Z stack-protector=strong} & Protección fuerte contra desbordamientos. \\
\texttt{-Z stack-protector=basic} & Protección básica de pila. \\
\texttt{-Z stack-protector=none} & Sin protección de pila. \\
\texttt{-C link-arg=-Wl,-z,noexecstack} & Memoria no ejecutable. \\
\texttt{-C link-arg=-Wl,-z,execstack} & Memoria ejecutable. \\
\texttt{-C relocation-model=pic -C link-arg=-pie} & \emph{PIE} habilitado. \\
\texttt{-C relocation-model=static -C link-arg=-no-pie} & \emph{PIE} deshabilitado. \\
\hline
\end{tabularx}
\caption{\emph{Flags} de seguridad en \emph{Rustc} [19]}
\label{tab:rust_seguridad}
\end{table}

\section{Análisis de resultados}
Para los resultados se usan medias de varias pruebas realizadas con los distintos programas.

\subsection{Tiempo de ejecución}

\subsubsection{Impacto de las Optimizaciones (\texttt{-O0} a \texttt{-Ofast})}
En la Tabla \ref{tab:optimizaciones_por_compilador} se presentan los tiempos de ejecución de los programas sin aplicar ninguna medida de protección. Esto permite establecer una línea base para cada nivel de optimización, con el objetivo de compararlos posteriormente con los tiempos obtenidos al activar medidas de seguridad y así evaluar el impacto que estas tienen en cada nivel de optimización. Además, se busca analizar la mejora relativa entre los distintos niveles de optimización para entender mejor cómo influyen dichas medidas en relación con el comportamiento del archivo base.


\begin{table}[H]
\centering
\footnotesize
\begin{tabular}{|c|c|c|c|}
\hline
\textbf{Nivel} & \textbf{\emph{Clang}} & \textbf{\emph{GCC}} & \textbf{\emph{Rustc}} \\
\hline
\texttt{-O0}    & 1.34 ms         & 1.40 ms        & 1.66 ms        \\
\texttt{-O1}    & 1.33 ms (-0.7\%) & 1.74 ms (+24\%) & 1.35 ms (-19\%) \\
\texttt{-O2}    & 1.55 ms (+16\%)  & 1.08 ms (-23\%) & 2.32 ms (+40\%) \\
\texttt{-O3}    & 1.31 ms (-2\%)   & 1.23 ms (-12\%) & 1.52 ms (-8\%)  \\
\texttt{-Os}    & 1.74 ms (+30\%)  & 1.12 ms (-20\%) & 1.35 ms (-19\%) \\
\texttt{-Oz}    & 2.00 ms (+49\%)  & N/A         & 1.53 ms (-8\%)  \\
\texttt{-Ofast} & 1.34 ms (=)      & N/A         & N/A          \\
\hline
\end{tabular}
\caption{Impacto de las optimizaciones por compilador}
\label{tab:optimizaciones_por_compilador}
\end{table}

En el nivel sin optimización (\texttt{-O0}), \emph{Clang} registra un tiempo de 1.34 ms, \emph{GCC} es ligeramente más lento con 1.40 ms, y \emph{Rustc} resulta aproximadamente un 24\% más lento que \emph{Clang} debido a las verificaciones de seguridad integradas. Al aplicar la optimización \texttt{-O1}, \emph{Clang} mantiene un rendimiento estable con una ligera mejora, \emph{GCC} sufre un retroceso notable y \emph{Rustc} mejora significativamente, superando a \emph{GCC}, lo que indica una mejor optimización en este nivel. Con \texttt{-O2}, \emph{GCC} destaca con la mejor mejora en velocidad, mientras que \emph{Clang} empeora y \emph{Rustc} se penaliza por los checks adicionales que realiza. En \texttt{-O3}, \emph{Clang} y \emph{GCC} alcanzan tiempos cercanos, con mejoras respecto a niveles previos, mientras que \emph{Rustc} se recupera, aunque sigue siendo el más lento. Al usar la optimización para tamaño \texttt{-Os}, \emph{GCC} ofrece un buen equilibrio con una mejora considerable, mientras \emph{Clang} empeora y \emph{Rustc} mantiene un rendimiento aceptable. En \texttt{-Oz}, \emph{Clang} sacrifica velocidad por un tamaño más reducido, \emph{Rustc} mantiene un rendimiento similar a \texttt{-O3}, y no hay datos disponibles para \emph{GCC}. Finalmente, \texttt{-Ofast}, disponible solo para \emph{Clang}, mantiene un rendimiento similar al nivel sin optimización pero omite comprobaciones de seguridad para maximizar la velocidad. En resumen, \emph{GCC} presenta el mejor rendimiento en niveles como \texttt{-O2} y \texttt{-Os}, \emph{Clang} es consistente y rápido en general, especialmente con \texttt{-Ofast}, y \emph{Rustc}, aunque generalmente más lento, mejora significativamente en ciertos niveles pese a sus medidas de seguridad.


\subsubsection{Impacto de las medidas de seguridad}

El objetivo es calcular el impacto exacto que implica activar protecciones de seguridad en aquellas combinaciones donde se observaron los peores resultados de rendimiento. Para ello, se evalúa el impacto de las \emph{flags} de seguridad en cada compilador utilizando su nivel de optimización óptimo, determinado según los datos de rendimiento previamente analizados: en el caso de \emph{Clang}, se selecciona \texttt{-O3} por su rendimiento consistente; para \emph{GCC}, se opta por \texttt{-O3}, ya que fue el mejor tiempo en las pruebas; y para \emph{Rustc}, se utiliza también \texttt{-O3}, al ofrecer un buen equilibrio entre seguridad y rendimiento. Los datos extraídos de las pruebas pueden consultarse en las Figuras \ref{fig:Gráficas_gráficas_por_optimización__o2_time_ms} y \ref{fig:Gráficas_gráficas_por_optimización__o3_time_ms}.

\begin{table}[H]
\centering
\footnotesize
\begin{tabular}{|l|c|c|c|}
\hline
\emph{\textbf{Flag}} & \textbf{\emph{Clang} (\texttt{-O3})} & \textbf{\emph{GCC} (\texttt{-O3})} & \textbf{\emph{Rustc} (\texttt{-O3})}\\
\hline
\texttt{-fstack-protector-strong} & 1.31 $\rightarrow$ 1.35 ms & 1.08 $\rightarrow$ 1.65 ms & 1.52 $\rightarrow$ 1.39 ms\\
\texttt{-D\_FORTIFY\_SOURCE=2} (C/C++) & 1.31 $\rightarrow$ 1.23 ms & 1.08 $\rightarrow$ 1.23 ms & N/A\\
\hline
\multicolumn{4}{|l|}{\textbf{Sanitizadores:}} \\
\hline
\texttt{-fsanitize=address} & 1.31 $\rightarrow$ 1.46 ms & 1.08 $\rightarrow$ 1.78 ms & 1.52 $\rightarrow$ 2.40 ms\\
\texttt{-Z sanitizer=memory} (Rust) & N/A & N/A & 1.52 $\rightarrow$ 2.71 ms\\
\hline
\end{tabular}
\caption{Casos con mayor impacto de flags de seguridad en niveles óptimos}
\label{tab:overhead_optimos}
\end{table}

Tras las pruebas y observando los peores resultados, se han resumido en la Tabla \ref{tab:overhead_optimos} los casos donde las \emph{flags} de seguridad causan el mayor impacto en el rendimiento para cada compilador en su nivel óptimo de optimización (\texttt{-O3}). Como se puede observar, \emph{GCC} presenta el mayor incremento relativo de tiempo de ejecución al activar ciertas protecciones, especialmente con el uso del \texttt{-fstack-protector-strong}. Por otro lado, \emph{Rustc} mantiene un coste relativamente bajo al emplear las \emph{flags} estándar, como las verificaciones de desbordamiento (\emph{overflow checks}), aunque los sanitizadores como el \texttt{-Z sanitizer=memory} aumentan significativamente el tiempo de ejecución. \emph{Clang} muestra un comportamiento intermedio, con un impacto moderado al activar estas protecciones.

\subsubsection{Comparativa final: seguridad vs. rendimiento}

En la Tabla \ref{tab:seguridad_vs_rendimiento} se ha elegido un nivel típico para producción como \texttt{-O3}
y se han comparado los tiempos de ejecución de cada compilador sin \emph{flags} de seguridad junto al impacto
que suponen aquellas comunes en este ámbito como \texttt{-fstack-protector-strong} y \texttt{-D\_FORTIFY\_SOURCE=2}.

\begin{table}[H]
\centering
\footnotesize
\begin{tabular}{|c|c|c|l|}
\hline
\textbf{Compilador} & \textbf{Tiempo (ms)} & \textbf{Impacto seguridad} & \textbf{Ventaja principal} \\
\hline
\emph{Clang} & 1.55 & -1.53\% (con \emph{flags}) & Optimización estable. \\
\emph{GCC}     & 1.08 & +33.34\% (con \emph{flags}) & Mejor rendimiento en \texttt{-O3}. \\
\emph{Rustc}   & 2.32 & -8.55\% (con \emph{flags}) & Seguridad sin configurar. \\
\hline
\end{tabular}
\caption{Comparativa de seguridad vs rendimiento}
\label{tab:seguridad_vs_rendimiento}
\end{table}

\subsubsection{Conclusiones}

En código con alta demanda de rendimiento la opción recomendada es\\
\emph{Clang}\ \texttt{-O3}\ \texttt{-D\_FORTIFY\_SOURCE=2\ -fstack-protector-strong}.
Aunque sufre un impacto combinado que reduce el tiempo en un 1.53\,\%,
siendo el más rápido. Esta opción se considera la ideal para entornos donde se
prioriza el rendimiento pero se requiere cierta protección básica.
En \emph{GCC}\ se debe evitar \texttt{-O3}\ \texttt{-D\_FORTIFY\_SOURCE=2\ -fstack-protector-strong}
ya que su tiempo aumenta en un 33.34.

En código donde la seguridad es crítica la mejor opción es
\emph{Rustc}\ \texttt{-O3} ya que mantiene un excelente equilibrio, posee
seguridad por defecto, y puede llegar a reducirse en 8.55\,\% el tiempo al
sumar protecciones. Además, no requiere configuración adicional para la
mayoría de los checks (desbordamiento, acceso inválido, etcétera).

En desarrollo, pruebas y depuración para C/C++, es muy recomendable el uso de
\texttt{-fsanitize=address} junto con los niveles de optimización \texttt{-O1} o \texttt{-O3} para
detectar errores de memoria. Aunque presenta un impacto elevado
($\sim$+64.81\,\%), resulta sumamente útil para capturar fallos complejos y difíciles de reproducir.
En el caso de Rust, la opción
\texttt{-Z\ sanitizer=memory} permite detectar accesos inválidos y condiciones de carrera con un impacto
similar en el rendimiento ($\sim$+78.28\,\%), siendo una herramienta eficaz en entornos de prueba. Estas variaciones
pueden ser observadas en la Figura \ref{fig:Gráficas_gráficas_por_optimización_porcentual__o3_time_ms}.

Los sanitizadores deben evitarse en producción, ya que pueden multiplicar el tiempo de ejecución hasta por 1.8×. 
Rust se destaca por ofrecer protecciones de memoria integradas con un impacto mínimo en el rendimiento, 
lo que lo convierte en una opción sólida cuando la seguridad no es negociable. \emph{GCC}, por otro lado, alcanza 
el mejor rendimiento bruto incluso al aplicar medidas de seguridad, aunque requiere ajustes más precisos y presenta 
un comportamiento menos predecible en niveles de optimización intermedios.

\subsection{Uso de memoria}

\subsection*{Efecto de las optimizaciones}

Las optimizaciones afectan el uso de memoria de manera diferente según el compilador y el nivel aplicado. A continuación, se detalla el consumo exacto en KB en las Tablas \ref{tab:memoria_clang}, \ref{tab:memoria_gcc} y \ref{tab:memoria_rustc} para \emph{Clang}, \emph{GCC} y \emph{Rustc}, respectivamente.

\begin{table}[H]
\centering
\footnotesize
\begin{tabular}{|l|c|c|}
\hline
\textbf{Nivel de optimización} & \textbf{Memoria (KB)} & \textbf{Variación respecto a \texttt{-O0}} \\
\hline
\texttt{-O0} (sin optimización) & 3270 & — \\
\texttt{-O1} & 3446 & \textbf{+5.4\%} \\
\texttt{-O2} & 3489 & \textbf{+6.7\%} \\
\texttt{-O3} & 3433 & \textbf{+5.0\%} \\
\texttt{-Os} (optimizado para tamaño) & 3416 & \textbf{+4.5\%} \\
\texttt{-Oz} (máxima reducción de tamaño) & \textbf{3013} & \textbf{-7.9\%}\\
\texttt{-Ofast} & 3448 & \textbf{+5.4\%} \\
\hline
\end{tabular}
\caption{Uso de memoria en \emph{Clang} según nivel de optimización}
\label{tab:memoria_clang}
\end{table}


\begin{table}[H]
\centering
\footnotesize
\begin{tabular}{|l|c|c|}
\hline
\textbf{Nivel de optimización} & \textbf{Memoria (KB)} & \textbf{Variación respecto a \texttt{-O0}} \\
\hline
\texttt{-O0} & 3052 & — \\
\texttt{-O1} & 3051 & \textbf{-0.03\%}\\
\texttt{-O2} & 3052 & \textbf{0\%} \\
\texttt{-O3} & 3459 & \textbf{+13.3\%}\\
\texttt{-Os} & \textbf{3038} & \textbf{-0.5\%}\\
\hline
\end{tabular}
\caption{Uso de memoria en \emph{GCC} según nivel de optimización}
\label{tab:memoria_gcc}
\end{table}


\begin{table}[H]
\centering
\footnotesize
\begin{tabular}{|l|c|c|}
\hline
\textbf{Nivel de optimización} & \textbf{Memoria (KB)} & \textbf{Variación respecto a \texttt{-O0}} \\
\hline
\texttt{-O0} & \textbf{1919} & — \\
\texttt{-O1} & 1996 & \textbf{+4.0\%} \\
\texttt{-O2} & 2004 & \textbf{+4.4\%} \\
\texttt{-O3} & 1966 & \textbf{+2.4\%} \\
\texttt{-Os} & 2014 & \textbf{+4.9\%} \\
\texttt{-Oz} & 1976 & \textbf{+3.0\%} \\
\hline
\end{tabular}
\caption{Uso de memoria en \emph{Rustc} según nivel de optimización}
\label{tab:memoria_rustc}
\end{table}

En general, las optimizaciones tienden a aumentar ligeramente el uso de memoria en \emph{Clang}, salvo en el caso de \texttt{-Oz}, que consigue una reducción significativa de aproximadamente un 8\%. Por su parte, \emph{GCC} mantiene un consumo estable en los niveles \texttt{-O0}, \texttt{-O1} y \texttt{-O2}; sin embargo, \texttt{-O3} provoca un aumento del 13\% en el uso de memoria, mientras que \texttt{-Os} logra una ligera reducción. Finalmente, \emph{Rustc} destaca por tener, en general, el menor consumo de memoria, con un incremento moderado y controlado de alrededor del 2 a 5\% debido a las optimizaciones, sin presentar cambios bruscos.


En base a los resultados obtenidos, \texttt{-Oz} (\emph{Clang}) y \texttt{-Os} (\emph{GCC}) son los únicos niveles que reducen memoria respecto a \texttt{-O0}. \texttt{-O3} es el que más aumenta memoria en \emph{GCC} (13\%), pero en \emph{Clang} y \emph{Rustc} el impacto es menor ($\sim$5\%). El compilador más consistente y eficiente es \emph{Rustc}, que siempre está por debajo de 2100 KB.

\subsection*{Efecto de las medidas de seguridad}

Algunas protecciones incrementan significativamente el uso de memoria. Los casos más notables incluyen los sanitizadores, que son los mayores consumidores. Por ejemplo, \emph{AddressSanitizer} puede aumentar la memoria en \emph{Clang} desde aproximadamente 3400 KB hasta $\sim$12000--14000 KB (\texttimes3--4). En \emph{Rustc}, el aumento es de $\sim$2000 KB a $\sim$8000--10000 KB (\texttimes4--5).

Respecto a las protecciones de pila, usar \texttt{-fstack-protector-strong} aumenta la memoria entre 5 y 10\% en \emph{Clang} y \emph{GCC}. Con \texttt{-fstack-protector-all} el aumento puede llegar al 15--20\%, ambos datos son visibles en las Ficguras \ref{fig:Gráficas_gráficas_por_medida_seguridad_porcentual_stack_protection_strong_memory_usage_kb_porcentaje} y \ref{fig:Gráficas_gráficas_por_medida_seguridad_porcentual_stack_protection_all_memory_usage_kb_porcentaje} respectivamente.

Otras protecciones como PIE tienen un impacto menor (2--5\%) como se puede ver en la Figura \ref{fig:Gráficas_gráficas_por_medida_seguridad_porcentual_pie_enabled_memory_usage_kb_porcentaje}. La opción \texttt{-flto} puede incluso reducir ligeramente el uso de memoria (3--5\%).

Rust tiende a manejar mejor ciertas protecciones como el stack protector (aumentos de 5--10\%), mientras que en C++ se elevan más (15--20\%). \emph{Clang} y \emph{GCC} muestran comportamientos similares, aunque \emph{GCC} suele ser ligeramente más eficiente al aplicar sanitizers.

\subsection*{Comparación entre compiladores}

Para estas comparaciones se han utilizado el consumo obtenido en \texttt{-O0} y \texttt{-O3} sin \emph{flags} de seguridad en la Tabla \ref{tab:compiladores_sin_prote} para ver como varía en función del grado de optimización, también se ha hecho la comparativa con algunas protecciones intensivas como \emph{AddressSanitizer}, mostrada en la Tabla \ref{tab:memoria_intensiva}.

\begin{table}[H]
\centering
\footnotesize
\begin{tabular}{|c|c|c|}
\hline
\textbf{Compilador} & \textbf{\texttt{-O0} (KB)} & \textbf{\texttt{-O3} (KB)} \\
\hline
\emph{Clang} & 3270 & 3433 \\
\emph{GCC} & 3052 & 3459 \\
\emph{Rustc} & \textbf{1919} & \textbf{1966} \\
\hline
\end{tabular}
\caption{Consumo en memoria comparado entre compiladores sin protecciones}
\label{tab:compiladores_sin_prote}
\end{table}

\begin{table}[H]
\centering
\footnotesize
\begin{tabular}{|c|c|c|}
\hline
\textbf{Compilador} & \textbf{Memoria base (KB)} & \textbf{Con ASan (KB)} \\
\hline
\emph{GCC} & $\sim$3050 & $\sim$11000--13000 \\
\emph{Clang} & $\sim$3400 & $\sim$12000--14000 \\
\emph{Rustc} & $\sim$2000 & $\sim$8000--10000 \\
\hline
\end{tabular}
\caption{Impacto de \emph{AddressSanitizer} en el uso de memoria}
\label{tab:memoria_intensiva}
\end{table}


\textbf{Optimización para tamaño (\texttt{-Os} / \texttt{-Oz}):}

En cuanto al uso de memoria, \emph{Clang} con la opción \texttt{-Oz} obtuvo un consumo de 3013 KB, siendo este el mejor resultado entre los compiladores de C++. Por su parte \emph{GCC}, usando \texttt{-Os}, presentó un uso ligeramente superior, con 3038 KB. Finalmente, \emph{Rustc} con \texttt{-Oz} demostró ser el más eficiente en términos absolutos, con un consumo de memoria de tan solo 1976 KB como se puede ver en las Figuras \ref{fig:Gráficas_gráficas_por_optimización__os_memory_usage_kb} y \ref{fig:Gráficas_gráficas_por_optimización__oz_memory_usage_kb}.

\subsection*{Conclusión final}

Rust es el más eficiente en cuanto al uso de memoria en todos los escenarios analizados, mostrando un consumo estable incluso al aplicar distintas medidas de seguridad. Por su parte, \emph{Clang} y \emph{GCC} presentan comportamientos comparables, aunque \emph{Clang} logra un mejor rendimiento cuando se compila con la opción \texttt{-Oz} orientada a reducir el tamaño.

Las protecciones que generan un mayor impacto en el uso de memoria son los sanitizadores, ya que pueden multiplicar el consumo por un factor de 3 a 5 veces. En cuanto a las optimizaciones, la opción \texttt{-O3} no siempre mejora el uso de memoria: en \emph{GCC}, por ejemplo, puede incrementar el consumo en un 13\%, mientras que en \emph{Rustc} su efecto es mucho más limitado, con un aumento cercano al 2\%.

\subsection{Tamaño del binario}

\subsubsection{\texorpdfstring{\textbf{Optimización}}{Optimización}}

Las opciones de compilación como \texttt{-Os} y \texttt{-Oz} están
diseñadas para reducir el tamaño de los archivos binarios que se generan
al compilar un programa.
Según los datos analizados, estas opciones logran binarios más pequeños
que los niveles de optimización estándar como \texttt{-O0}, \texttt{-O1}
o \texttt{-O2}. Entre ellas, \texttt{-Oz} es más efectiva en esta tarea
que \texttt{-Os}, ya que logra tamaños incluso menores en algunos casos,
con diferencias que pueden llegar hasta los 6000 KB, como se puede ver
en las Figuras \ref{fig:Gráficas_gráficas_por_optimización__os_file_size_kb} y 
\ref{fig:Gráficas_gráficas_por_optimización__oz_file_size_kb}.

Por otro lado, la opción \texttt{-Ofast} (véase la Figura 
\ref{fig:Gráficas_gráficas_por_optimización__ofast_file_size_kb}), aunque no está disponible en
todos los compiladores, produce binarios más pequeños que los generados
con \texttt{-O3}, pero siguen siendo más grandes que los de \texttt{-Os}
o \texttt{-Oz}, ya que está enfocada en la velocidad de ejecución y no
tanto en el tamaño, aunque algunas optimizaciones ayudan a eliminar
código innecesario.

En cuanto a los niveles más agresivos como \texttt{-O3} y
\texttt{-Ofast}, se observa que generan binarios considerablemente más
grandes, especialmente con \emph{Clang} y \emph{GCC}, debido a
técnicas como la inserción masiva de funciones (\emph{inlining}) y la
vectorización del código.

Aunque \texttt{-Ofast} puede resultar en un tamaño ligeramente menor que
\texttt{-O3}, sigue teniendo un tamaño mayor que las opciones enfocadas
en reducir el tamaño como \texttt{-Os}.

\subsubsection{\texorpdfstring{\textbf{Seguridad}}{Seguridad}}

Cuando se activan ciertas protecciones en la compilación de programas,
el tamaño de los binarios puede aumentar notablemente, tal como muestran
las Figuras \ref{fig:Gráficas_gráficas_por_optimización_porcentual__os_file_size_kb} y
\ref{fig:Gráficas_gráficas_por_optimización_porcentual__oz_file_size_kb}.

Por ejemplo, los sanitizadores como \emph{AddressSanitizer} o \emph{MemorySanitizer}
pueden hacer que el binario final sea entre dos y tres veces más grande
debido a que se añade código adicional para revisar que no haya errores
mientras el programa se está ejecutando.

Otras protecciones como los \emph{stack protectors} también incrementan
el tamaño, aunque en menor medida. Además, mecanismos como RELRO
y {PIE} pueden sumar entre un 10\% y 20\% de tamaño adicional,
según el compilador.

El problema se intensifica cuando se combinan varias protecciones
simultáneamente. Por ejemplo, usar sanitizadores junto con protecciones
de pila y PIE puede llevar a que el binario en \emph{Clang}
supere los 5000 KB, comparado con unos 1000 KB que tiene en una
combinación básica sin \emph{flags} de seguridad.

\subsubsection{\texorpdfstring{\textbf{Comparación entre
compiladores}}{Comparación entre Compiladores}}

En general, los programas compilados con Rust suelen ocupar menos
espacio que los generados con compiladores de C++ como \emph{Clang} o
\emph{GCC}, incluso cuando se usan combinaciones similares.

Por ejemplo, la combinación por defecto en Rust produce binarios más
pequeños que la combinación por defecto en C++. Entre los compiladores
de C++, \emph{Clang} tiende a generar archivos ligeramente más
pequeños que \emph{GCC}, sobre todo cuando se usan opciones como
\texttt{-Os} o \texttt{-Oz}, pensadas para reducir el tamaño.

Al aplicar protecciones adicionales como sanitizadores el aumento de
tamaño es más notable en C++, usar el sanitizador de \emph{Address}
puede duplicar el tamaño del binario, mientras que en Rust el aumento es
más moderado, alrededor de un 50\%. Y se observa que el peor es \emph{Clang}
con hasta un 10000\% de tamaño superior, lo cual se sale mucho de la media.

En cuanto a otras protecciones como RELRO o protecciones de pila,
su efecto sobre el tamaño es similar en todos los compiladores, sin
embargo Rust parece integrar estas medidas de forma más eficiente, ya
que los binarios que genera son menores en tamaño que los de C++.

\subsubsection{\texorpdfstring{\textbf{Recomendaciones}}{Recomendaciones}}

Si lo que se quiere es reducir al máximo el tamaño de un programa, lo
más efectivo es compilarlo con la opción \texttt{-Oz} o \texttt{-Os}. Estas opciones le indican al
compilador que priorice el tamaño del ejecutable por encima del
rendimiento, lo que ayuda a que el resultado final ocupe lo menos
posible. También conviene desactivar herramientas de depuración o
protecciones adicionales que no sean estrictamente necesarias, ya que
suelen aumentar el tamaño del binario.

Ahora bien, si se busca un equilibrio entre tamaño y seguridad, una
buena opción es usar \texttt{-Os} combinado con medidas de protección
como \texttt{-fstack-protector-strong}. A esto se le pueden añadir opciones al enlazador como
\texttt{-Wl,-z,relro,-z,now}, que activan protecciones en tiempo de
ejecución para dificultar la modificación de partes sensibles del
programa.

Por último, no es recomendable usar muchas medidas pesadas al mismo
tiempo como los sanitizadores, PIE o protecciones muy
estrictas para la pila a no ser que sea realmente necesario. Este tipo
de combinaciones puede hacer que el programa aumente bastante de tamaño
y también que tarde más en ejecutarse.

\subsection{Funciones fortificadas}

Una de las principales herramientas para aumentar la seguridad del
código en C y C++ es la opción \texttt{FORTIFY\_SOURCE} [13]. Estas opciones
provocan que el compilador inserte
comprobaciones adicionales en ciertas funciones estándar que podrían ser
un vector de ataque, como los errores de desbordamiento de búfer. En
cambio, el compilador de Rust no admite esta opción, lo que se refleja
en los archivos analizados como ``N/A''.

Los sanitizadores, pese a que su propósito principal no es la
fortificación de funciones, contribuyen a la seguridad detectando
errores de memoria y comportamiento indefinido durante la ejecución del
programa.

La mejor combinación de medidas observada para maximizar la
fortificación es usar \texttt{-D\_FORTIFY\_SOURCE=2} junto con
sanitizadores a partir del nivel \texttt{-O1}, como se muestra en la Figura 
\ref{fig:Gráficas_gráficas_por_optimización__o1_file_size_kb}. 
Este tipo de combinaciones son las que provocan un mayor
porcentaje de funciones protegidas frente a errores, en comparación con
usar \texttt{FORTIFY\_SOURCE=1} por sí solo.
Pese a que Rust no implementa \texttt{FORTIFY\_SOURCE}, sus
sanitizadores pueden añadir una capa adicional de seguridad,
complementando el modelo de seguridad del lenguaje.
Por parte de \emph{GCC} y \emph{Clang}, soportan tanto
\texttt{FORTIFY\_SOURCE} como los sanitizadores. No se observan
diferencias notables en la cantidad de funciones fortificadas cuando se
usan las mismas opciones.

Para finalizar es necesario comentar que aunque Rust no posea soporte para
\texttt{FORTIFY\_SOURCE}, esto no implica menor seguridad.
El lenguaje está diseñado para prevenir errores comunes de memoria
mediante su sistema de tipos y el modelo de \emph{ownership} (gestión de la memoria asignando
un único propietario a cada valor y liberándola automáticamente al finalizar su uso), lo que
reduce la necesidad de fortificación adicional.

Por otro lado, las opciones de optimización como \texttt{-Ofast},
\texttt{-Os} o \texttt{-Oz} no influyen directamente en la
fortificación.

